\subsection{Purpose}

The purpose of this section is to establish minimum requirements for collecting, protecting, retaining, reviewing, and responding to logs and monitoring information from systems, applications, network devices, and security controls.

These requirements are intended to:
\begin{itemize}
    \item Detect unauthorized access, misuse, security incidents, and control failures
    \item Support investigation, incident response, troubleshooting, and recovery
    \item Provide accountability for administrative, privileged, and security-sensitive actions
    \item Ensure that important systems and services generate appropriate logs
    \item Protect logs from unauthorized access, alteration, deletion, or premature destruction
    \item Establish appropriate log retention, review, and disposal practices
    \item Provide timely notification of significant security and operational events
    \item Support the assessment and improvement of security, availability, and compliance controls
\end{itemize}

Systems, applications, network devices, and security controls \must{} generate sufficient logs and monitoring information to support security detection, troubleshooting, incident response, availability management, and compliance requirements.

Logging and monitoring requirements \must{} be based on the importance of the system, the sensitivity of the information, the likelihood of threats, and applicable legal, contractual, or business requirements.

\subsection{Logging Requirements}

Critical systems and security-relevant services \must{} produce logs for significant security, administrative, operational, and access events.

Where supported, systems \should{} log:
\begin{itemize}
    \item Successful and unsuccessful authentication attempts
    \item Account creation, modification, disabling, and deletion
    \item Changes to privileges, groups, roles, and access permissions
    \item Administrative and privileged actions
    \item Access to sensitive or restricted information
    \item Security-control changes
    \item Configuration changes
    \item Software installation, removal, and updates
    \item System startup, shutdown, and service failures
    \item Malware, intrusion, and other security alerts
    \item Data export, deletion, or transfer events
    \item Backup, restoration, and recovery events
    \item Network connections and remote-access sessions
\end{itemize}

Logging \must{} be enabled before systems are placed into production unless a documented exception has been approved.

\subsection{Log Content}

Logs \should{} contain enough information to support investigation and correlation. Where applicable, each event \should{} include:
\begin{itemize}
    \item Date and time
    \item Source system or device
    \item User, account, process, or service involved
    \item Source and destination address or resource
    \item Event type and outcome
    \item Affected system, application, account, or data
    \item Relevant transaction, session, or correlation identifier
    \item Description of the action or error
\end{itemize}

Systems \must{} use synchronized and reasonably accurate time sources. Time synchronization \should{} be centrally managed where practical.

Logs \mustnot{} contain passwords, private keys, authentication tokens, recovery codes, or other sensitive credential values. Sensitive information that is not required for logging \should{} be masked, removed, or replaced with a non-sensitive identifier.

\subsection{Centralized Log Collection}

Logs from critical systems and security devices \should{} be sent to a centralized logging or security-monitoring system where technically feasible.

Centralized logging \should{} be used to:
\begin{itemize}
    \item Reduce the risk of log loss or unauthorized alteration
    \item Correlate events across systems
    \item Support detection of suspicious activity
    \item Provide a consistent search and investigation capability
    \item Alert responsible personnel about important events
    \item Support incident response and forensic analysis
\end{itemize}

Log transmission and storage \must{} be protected against unauthorized access and alteration. Logs transmitted across untrusted networks \should{} use encrypted connections.

Systems \should{} continue producing local logs when centralized collection is unavailable. Collection failures \should{} generate an alert or other notification for responsible personnel.

\subsection{Log Protection and Access}

Access to logs \must{} be limited to authorized personnel with a business or operational need. Administrative access to logging systems \must{} use individual accounts where supported and \should{} require MFA.

Log access:
\begin{itemize}
    \item \must{} be restricted according to least privilege
    \item \must{} be monitored where feasible
    \item \must{} not allow ordinary users to alter security logs
    \item \should{} be separated from the administration of the systems being monitored
    \item \must{} be removed when the user's role or responsibilities no longer require it
\end{itemize}

Logs \must{} be protected from unauthorized deletion, modification, or premature destruction. Where appropriate, logs \should{} be stored using access controls, integrity protection, write-once storage, or other tamper-resistant mechanisms.

\subsection{Log Retention}

Logs \must{} be retained for a period appropriate to the system's importance, the sensitivity of the information, investigation requirements, and applicable legal, contractual, or business requirements.

The organization \should{} define retention periods for:
\begin{itemize}
    \item Authentication and access logs
    \item Administrative activity logs
    \item Security alerts
    \item Network and firewall logs
    \item Application and database logs
    \item Backup and restoration logs
    \item Logs associated with security incidents
\end{itemize}

Logs associated with an investigation, legal hold, audit, or unresolved incident \must{} be preserved until the responsible owner authorizes their release or disposal.

Log retention and disposal \must{} consider the sensitivity of the information contained in the logs. Logs \must{} be securely deleted or destroyed when the retention period expires, unless they are required for an active investigation or other approved purpose.

\subsection{Monitoring and Alerting}

Security and operational monitoring \must{} be implemented for systems and services based on their criticality and risk.

Monitoring \should{} identify:
\begin{itemize}
    \item Repeated authentication failures
    \item Authentication from unusual locations, devices, or times
    \item Privilege escalation or unexpected administrative activity
    \item Creation of unauthorized accounts or access paths
    \item Changes to security controls or firewall rules
    \item Malware, intrusion, or suspicious endpoint activity
    \item Unusual data access, transfer, or deletion
    \item Unexpected service failures or system shutdowns
    \item Resource exhaustion or capacity problems
    \item Backup failures
    \item Repeated network connection or access-denied events
\end{itemize}

Alerts \must{} be assigned an appropriate severity and routed to personnel who are responsible for reviewing and responding to them. Alert thresholds \should{} be adjusted to reduce unnecessary alerts while maintaining the ability to detect significant activity.

High-severity alerts \must{} be reviewed promptly. Monitoring personnel \must{} escalate suspected compromise, data loss, or significant control failure according to the incident-response procedure.

\subsection{Review and Investigation}

Logs and alerts \must{} be reviewed at a frequency appropriate to the system's risk and criticality. Critical security alerts \should{} be reviewed as soon as practicable.

Log reviews \should{} determine whether:
\begin{itemize}
    \item Suspicious or unauthorized activity occurred
    \item Administrative actions were expected and approved
    \item Security controls are operating effectively
    \item Systems are generating the required logs
    \item Time synchronization is functioning
    \item Log collection and storage are operating correctly
    \item Retention requirements are being met
\end{itemize}

Investigations \must{} preserve relevant logs and supporting evidence. Investigators \should{} document the event, affected systems, relevant time period, findings, actions taken, and required follow-up.

\subsection{Monitoring Privacy and Data Minimization}

Monitoring \must{} be limited to information necessary for security, operational, legal, or business purposes. Access to monitoring information \must{} be limited to authorized personnel.

Monitoring systems \should{} avoid collecting unnecessary personal information. Where personal information is collected, it \must{} be protected according to the applicable data-classification and retention requirements.

\subsection{Availability of Logging Systems}

Logging and monitoring systems supporting critical operations \should{} have appropriate backup, capacity, and recovery arrangements. Administrators \must{} monitor available storage and take action before logs are lost because of capacity limitations.

A failure of centralized logging \must{} be treated as a control failure and reported to the responsible administrator. The administrator \must{} restore logging as soon as practicable and determine whether any required events were lost during the failure.

\subsection{Review of Logging and Monitoring Controls}

Logging and monitoring configurations \must{} be reviewed periodically and after significant changes to systems, applications, networks, threats, or business requirements.

The review \should{} verify that:
\begin{itemize}
    \item Required systems and devices are generating logs
    \item Important events are being collected
    \item Logs contain sufficient information for investigation
    \item Time synchronization is functioning
    \item Alerts are routed to responsible personnel
    \item Log access is appropriately restricted
    \item Retention periods remain appropriate
    \item Storage capacity is sufficient
    \item Monitoring gaps and unresolved alerts are tracked
\end{itemize}

Systems that cannot meet required logging or monitoring controls \must{} be documented and approved as an exception. The exception \must{} identify the missing capability, associated risks, compensating controls, responsible owner, approval, and expiration or review date.